% Bagian: first-order-logic
% Bab: beyond
% Subbagian: other-logics

\documentclass[../../../include/open-logic-section]{subfiles}

\begin{document}

\olfileid[id]{fol}{byd}{oth}

\olsection{Logika-Logika Lain}

Seperti mungkin telah Anda pahami sekarang, merancang logika baru tidaklah
sulit. Anda juga dapat menciptakan sintaks sendiri, merancang sistem deduktif,
dan membentuk semantik yang menyertainya. Anda mungkin harus agak cerdik jika
ingin agar sistem !!{derivation} tersebut lengkap bagi semantiknya, dan
mungkin perlu upaya untuk meyakinkan khalayak luas bahwa logika Anda
benar-benar menarik. Namun sebagai imbalannya, Anda dapat menikmati
berjam-jam kegiatan yang menyenangkan dan bermanfaat dengan menjelajahi
sifat-sifat matematis dan komputasional logika Anda.

Beberapa dasawarsa terakhir telah menyaksikan ledakan logika formal yang
sungguh besar. Logika kabur dirancang untuk memodelkan penalaran tentang sifat
yang samar. Logika probabilistik dirancang untuk memodelkan penalaran tentang
ketidakpastian. Logika default dan logika nonmonoton dirancang untuk memodelkan
bentuk penalaran yang dapat dibatalkan, yakni inferensi ``masuk akal'' yang
kemudian dapat gugur ketika dihadapkan pada informasi baru. Terdapat logika
epistemik, yang dirancang untuk memodelkan penalaran tentang pengetahuan;
logika kausal, yang dirancang untuk memodelkan penalaran tentang hubungan
sebab-akibat; dan bahkan logika ``deontik'', yang dirancang untuk memodelkan
penalaran tentang kewajiban moral dan etis. Bergantung pada apakah motivasi
utama untuk memperkenalkan sistem-sistem ini bersifat filosofis, matematis,
atau komputasional, Anda dapat menjumpai makhluk-makhluk semacam itu dikaji di
bawah payung logika matematika, logika filosofis, kecerdasan buatan, ilmu
kognitif, atau bidang lain.

Daftarnya terus berlanjut, dan kemungkinannya tampak tanpa batas. Mungkin kita
tidak akan pernah mencapai impian Leibniz untuk mereduksi seluruh penalaran
manusia menjadi perhitungan---tetapi itu tidak dapat menghentikan kita untuk
mencobanya.

\end{document}
